%% Introduction du problème %%

Nous ne pouvons pas non plus ignorer que ces institutions patrimoniales conservent un nombre conséquent de données, renseignées pour la plupart manuellement. Il est donc hautement probable que des erreurs persistent parmi les données, parfois à cause d'un manque de référentiels et d'autres simplement à cause de l'erreur humaine. Autant sur certains champs comme le titre qu'au sein même de l'indexation, notamment de noms propres. Au cours de l'exploration des données d'une des bases de données du dépôt légal, nous avons ainsi relevé sans difficulté 3 émissions dont il existait deux entités différentes : \enquote{Arte Journal} et \enquote{ARTE journal}, \enquote{13h15 le dimanche} et \enquote{13h15 le dimanche.}, \enquote{Mardi cinéma, le dimanche aussi !} et \enquote{Mardi cinéma, le dimanche aussi. [Le mag].}. Ces erreurs peuvent empêcher certaines fiches de ressortir lors d'une recherche. La réalisation de statistiques peut également en être entachée. Ce sont les mêmes conséquences pour les indexations de noms propres. Pascal Lebeurier évoque dans son entretien les longues recherches nécessaires pour retrouver le nom de footballers mal indexés dans fiches anciennes\footnote{(Voir entretien de Pascal Lebeurier sur l'expertise du documentaliste ~\ref{ann:transcriptions},\nameref{ann:transcriptions}), p.~\pageref{ann:transcriptions}.}. Sans que d'erreurs manifestes ne soient commises les fiches peuvent également varier. Si deux notices de deux épisodes d'un même programme sont réalisés par deux personnes différentes, l'indexation appliquée peut varier. Nous ne pouvons pas parler \enquote{d'erreur} mais cela entraîne tout de même une variation dans les résultats selon la recherche qui est effectuée et tous documents voulus ne ressortirons pas. 